% ============================================================
% LAYOUT AND SPACING
% ============================================================

\setlength{\parskip}{3.5pt plus 1pt minus 0.5pt}
\setlength{\parindent}{0pt}

\widowpenalty=10000
\clubpenalty=10000
\brokenpenalty=4000
\predisplaypenalty=10000

\raggedcolumns
\newcounter{zsfFirstChap}\setcounter{zsfFirstChap}{1}

% Kompakte Spacing-Skala (4 Stufen)
\newlength{\ZSFspaceXS}          \setlength{\ZSFspaceXS}{1pt}
\newlength{\ZSFspaceS}           \setlength{\ZSFspaceS}{4pt}
\newlength{\ZSFspaceM}           \setlength{\ZSFspaceM}{7pt}
\newlength{\ZSFspaceL}           \setlength{\ZSFspaceL}{10pt}

% Abstand ober-/unterhalb von TextBlocks und Trennlinien
\newlength{\ZSFspaceOuter}        \setlength{\ZSFspaceOuter}{7pt}
\newlength{\ZSFspaceSep}          \setlength{\ZSFspaceSep}{3pt}

\setlength{\ZSFtextMinPad}{\ZSFspaceXS}

\setlength{\columnsep}{\ZSFspaceL} % Konsistent mit Skala
\setlength{\multicolsep}{\ZSFspaceS}
\setlength{\emergencystretch}{2em}
\newlength{\ZSFbreakThreshold}   \setlength{\ZSFbreakThreshold}{3\baselineskip} % Mindesthöhe für BreakIfNotEnoughSpace
\newlength{\ZSFpostChapterNudge} \setlength{\ZSFpostChapterNudge}{0pt}  % Korrektur nach Kapitelbalken (bei Bedarf anpassen)
\newlength{\ZSFtitleBarHeight}    \setlength{\ZSFtitleBarHeight}{9.2pt}  % feste Hoehe aller Titelbalken
\newlength{\ZSFsubsectionBarBeforeSkip}\setlength{\ZSFsubsectionBarBeforeSkip}{\ZSFspaceL}     % Standard-Abstand vor SubsectionBar
\newlength{\ZSFsubsectionAfterChapterGap}\setlength{\ZSFsubsectionAfterChapterGap}{\ZSFspaceXS} % Restabstand wenn SubsectionBar direkt nach ChapterBar steht
\newlength{\ZSFsubsectionAfterGap}       \setlength{\ZSFsubsectionAfterGap}{\ZSFspaceXS}        % Restabstand wenn Box direkt nach SubsectionBar steht
% Bewusste Feinjustierung zwischen XS und S: nicht in Skala, weil der Abstand
% i.d.R. durch \addvspace der SubsectionBar absorbiert wird und nur in den
% Sonderfällen ohne nachfolgende SubsectionBar visuell auftaucht.
\newlength{\ZSFchapterBarAfterSkip}      \setlength{\ZSFchapterBarAfterSkip}{2pt} % Standard-Abstand nach ChapterBar

% Tokens fuer statementbox / procedure (Aussagen & Verfahren)
\newlength{\ZSFstatementBarWidth} \setlength{\ZSFstatementBarWidth}{2pt}   % Breite des linken Akzentstreifens
\newlength{\ZSFstatementLeftPad}  \setlength{\ZSFstatementLeftPad}{\ZSFspaceM} % Abstand Streifen -> Text
\newlength{\ZSFprocStepGap}       \setlength{\ZSFprocStepGap}{\ZSFspaceXS} % Abstand Schritt-Titel -> Beschreibung
\newlength{\ZSFprocItemSep}       \setlength{\ZSFprocItemSep}{\ZSFspaceXS} % Abstand zwischen Schritten

% ============================================================
% TEXT-BOX-BINDUNGEN (Klassen für Abstände)
\newcommand{\ZSFRobustUnskip}{%
  \ifvmode
    \ifdim\lastskip>0pt\relax\vskip-\lastskip\fi
    \ifnum\lastpenalty=0\else\unpenalty\fi
    \ifdim\lastskip>0pt\relax\vskip-\lastskip\fi
  \else
    \ifdim\lastskip>0pt\relax\unskip\fi
    \ifnum\lastpenalty=0\else\unpenalty\fi
    \ifdim\lastskip>0pt\relax\unskip\fi
  \fi
}

% Text, der zur Box DARÜBER gehört (Post-Box Text)
\newcommand{\textNachBox}[1]{%
    \par
    \ZSFRobustUnskip % Löscht zu 100% verlässlich den Skip der Box (ohne manuelles Minus-Rechnen)
    \vspace{\ZSFspaceXS}% Unser exakter Soll-Abstand
    \begingroup
    \setlength{\parskip}{0pt}% Kein zusätzlicher Absatz-Abstand!
    \nopagebreak
    {\small\color{TextNachBoxColor} \noindent #1\par}%
    \endgroup
    \addvspace{\ZSFspaceOuter}% Abstand zum nächsten Element
}

% Text, der zur Box DARUNTER gehört (Pre-Box Text)
\newcommand{\textVorBox}[1]{%
    \par
    \addvspace{\ZSFspaceOuter}%
    \begingroup
    \setlength{\parskip}{0pt}% Kein zusätzlicher Absatz-Abstand!
    {\small\color{TextVorBoxColor} \noindent #1\par}%
    \endgroup
    \nopagebreak
    \vspace{\ZSFspaceXS}% Unser exakter Soll-Abstand
    \vspace{-\ZSFspaceS}% Kompensiert den kommenden Box-Abstand
    \nopagebreak% Kein Spaltenumbruch nach negativem Glue (verhindert Überlappung in multicols)
}

% ============================================================
% TEXT-FORMEL-BINDUNGEN (Innerhalb von formulabox)
% ============================================================

% Titel/Text vor einer Formel (innerhalb der formulabox)
\newcommand{\textVorFormel}[1]{%
    \par
    \begingroup
    \setlength{\parskip}{0pt}% Kein zusätzlicher Absatz-Abstand!
    \noindent\raggedright\textcolor{black!60}{\textbf{#1}}\par
    \endgroup
    \vspace{\ZSFspaceXS}% Unser exakter Soll-Abstand
    \centering
}

% Erklärungstext nach einer Formel (innerhalb der formulabox)
\newcommand{\textNachFormel}[1]{%
    \par
    \vspace{\ZSFspaceXS}% Unser exakter Soll-Abstand
    \begingroup
    \setlength{\parskip}{0pt}% Kein zusätzlicher Absatz-Abstand!
    \noindent\raggedright\textcolor{black!60}{\small #1}\par
    \endgroup
    \centering
}
